\documentclass[12pt]{book}

\usepackage{graphicx}
\usepackage{eso-pic}
\usepackage{tcolorbox}
\usepackage{polyglossia}
\setdefaultlanguage[indentfirst=false,forceheadingpunctuation=false]{russian}
\setotherlanguages{english}

\RequirePackage[a4paper, headsep = 0.5 \headsep, left=2.5cm, right=2.1cm, top=2cm, bottom=2.1cm]{geometry}

\setmainfont{URW Bookman}
\newfontfamily\cyrillicfont{URW Bookman}[Script=Cyrillic, Ligatures=TeX]
\setlength{\parindent}{20pt}

\usepackage{enumitem}
\setlist[itemize, 1]{label =\raisebox{-0.3\height}{\scalebox{1.8}{\textbullet}}}

\title{Феноменология чуждого}
\author{Бернхард Вальденфельс}
\date{2006}
\begin{document}

\AddToShipoutPictureBG*{\includegraphics[width=\paperwidth]{title.png}}

\null
    \thispagestyle{empty}%
    \addtocounter{page}{-1}%
    \newpage

\tableofcontents

\chapter{Грани чуждого}

Выделить чуждое как отдельную тему --- значит уже упустить его, начать со знакомого места и планировать к нему же вернуться. Обращение к чуждому наверняка приведет к переменам и даже катарсису, хотя в конечном счете привычность и возобладает и может даже углубиться. Чуждое не безобидно, и может отдалить нас от самих себя --- потому ему стараются сопротивляться, избегать его или ассимилировать. Эти старания оставляют субъекта наедине с самим собой и требуют выстраивать и оборонять порядок, исключающий проявление чуждого изнутри. Чуждое может пробудить любопытство и рассказать нам о нас самих, но как только оно подходит к тайне свободы и разума --- срабатывает сигнал ''Хаос!'', и свобода и разум остаются в своих правах. Иначе им пришлось бы отказаться от самих себя. Чуждость неизбежно ведет к враждебности, обе стороны все больше убеждаются в своей правоте. Мы становимся ближе и бдительней друг к другу. Встроенные в эмпирическую сеть защитные механизмы исходят из \textit{своего} и стремятся закрепиться в \textit{общем}. Чуждое предполагается внешним, ожидается, что оно будет идентифицировать себя и выглядеть незваным гостем. Его тогда можно будет подвергнуть оценке и суждению и поместить на повседневный моральный, политический, религиозный, культурный и интеллектуальный карантин.

Если же чуждое воспринимать как неизбежно неопределенное, как всегда сокрытое в \textit{своем}, где оно тревожит, соблазняет или пугает, ускользает от понимания, то опыт чуждого, значит, воздействует на \textit{наш} опыт, становясь отчуждением от опыта. Чуждость самореферентна и заразительна, ее последствия предшествуют всякой ее тематизации. Феноменология чуждого, предвосхищенная Гуссерлем в нескольких его новаторских концепциях (чуждый опыт, чуждое эго, чуждое тело и чуждый мир) обладает задатками науки, могущей говорить о чуждом во всей его полноте, всерьез принимать его радикальность, не довольствуясь сведением к его пониманию и практическому распознаванию. Феноменолог, всерьез принимающий эти требования ксенологии, наверняка будет вдохновляться работами Марселя Мосса, Георга Зиммеля, Вальтера Беньямина, Людвига Витгенштейна или Михаила Бахтина, не говоря уже о таких литераторах как Кафка, Музиль, Целан, Валери и Кальвино. В мире искусства другие порядки, и художественное чутье легко преодолевает границы самоуверенности и принципиальности, если вступает в контакт с чуждым.

Настоящее исследование не пытается охватить проблематику чуждого во всей ее пугающей полноте, а лишь обозначает некоторые основные мотивы, побудившие автора разработать феноменологию чуждого в соответствующей серии книг. Первый том --- \textit{Жало чуждого} вышел в 1990г.\footnote{Формальное изложение феноменологии чуждого можно найти в Главе 1 \textit{Топографии чуждого} (\textit{Topographie des Fremden}, 1997)}. Аспекты чуждого рассматриваются в главах 5-6 книги \textit{Рваные линии опыта} (\textit{Bruchlinien der Erfahrung}, 2002).} Мотивы здесь развивают тему, а не просто ретроактивно обосновывают обращение к ней. Опыт чуждого --- как и опыт удивления, ужаса, любви, ненависти или тех едва заметных несостыковок, через которые возникнет новое --- нельзя вполне описать или обосновать. Разделение на главы выполнено так, что каждая рассматривает конкретный программный компонент проблематики чуждого. Они выстроены в порядке этапов формирования феноменологии чуждого --- они потому вводные, поясняющие тексты и дополняют уже опубликованные обширные исследования. В последующих абзацах я свяжу эту работу с предыдущими, что избавит меня от необходимости отдельно на них ссылаться и позволит сосредоточиться на ключевых темах: порядок -- пафос -- ответ -- тело -- внимание -- межкультурность.

В главе 1 рассматривается \textit{Порядок в сумерках} (\textit{Ordnung im Zwielicht}, 1987). На границах всякого порядка чуждое предстает экстраординарным, не находящим места в порядке --- и все же оно, будучи исключаемым, не ничто. Поскольку оно исключается не вообще --- а из конкретного порядка, то оно нечто большее, чем просто сливающееся с фоном в неопределенности. И поскольку оно взывает к нашим чувствам, то оно не просто шум, от которого мы хотели бы укрыться. Положение, в котором мы оказываемся, обнаруживая несхватываемое в схватываемом, беспорядочное в упорядоченном, невидимое в видимом, тишину в озвучиваемом, запускает цепь исследований, сводящихся к вопросу: как мы можем взаимодействовать с чуждым без предварительной нейтрализации или отрицания его эффектов, его вызовов и требований? Как феноменолог, я всегда старался избегать того, в чем нетерпеливая аудитория любит обвинять феноменологов чуждого --- возведения его в ранг некоего морального или религиозного авторитета или ''совершенно иного''. Это было бы обожествлением Другого --- в чем часто обвиняют Левинаса --- и морализацией чуждого. Но навязать переживание столь же возможно, как и, согласно Канту, чувство. Не то чтобы восприятие чуждого было лишь субъективным ощущением, как в том убеждает модерн, просто в этом мире вещей гораздо больше, чем может охватить школа мысли с ее нормами, ценностями и интересами.

К переживаниям можно относиться с разной степенью осторожности: позволить им говорить самим за себя, использовать как информационный ресурс или пытаться исключить --- будь то из-за небрежности, властолюбия или безразличия. В главах 2 и 3 эти возможности рассмотрены с точки зрения ответовой этики --- этики, возникающей в ответе, отклике, реакции, и так оставающейся за рамками заповедей и запретов. До-предикативный опыт в гуссерлевском смысле здесь обращается в до-нормативный --- остающийся всегда позади, пока мы намечаем действия и строим планы, требующие моральной или правовой оценки. В главе 2 я буду ссылаться в основном на \textit{Рваные линии опыта} (\textit{Bruchlinien der Erfahrung}, 2002): анализ, проведенный в этом тексте, построен вокруг двойного события, состоящего из пафоса и последующей реакции. В начале есть то, что предстает как события --- происходящим с нами, в чем я или Другой можем участвовать --- а не как волящий и действующий субъект. Ключевое здесь --- временной сдвиг. Не нечто иное предшествует инициативе ''я'' --- эта инициатива предшествует самой себе, не возникая из неоткуда. Эти связи предстанут в новом свете, когда в главе 3, рассматривая \textit{Регистр Ответов} (\textit{Antwortregister}, 1994)\footnote{См. \textit{Феноменология внимания} (\textit{Phänomenologie der Aufmerksamkeit}), глава 3 \textit{Неисчислимое} (\textit{Unerzählbares}).}, я раскрываю основные черты ответового поведения и соответствующую логику ответов, возникающую в результате зова чуждого. Парадокс творческого ответа представляет особый случай: оказывается, мы можем воспринимать ''что'' ответа как нечто изобретенное, но не можем сделать то же с его ''на что''. Зов чуждого может быть столь же слабо связан с этическим фундаментализмом, сколь и с технологическим конструктивизмом.

Глава 4 посвящена телесному опыту, и не в коем случае не просто расширяет сферу чуждого. В узле самореференции и самоустранения, определяющих его способом существования, тело оказывается как бы эмблемой чуждого. Чужеродность начинается не с меня самого --- это мало что значило бы,--- так называемый субъект выпадает из центра, когда предстает как \textit{Воплощенное Я} (\textit{Leibliches Selbst})\footnote{Название лекций о теле, опубликованых в 2000г., в которых эта тема рассматривается во всей ее глубине.}. В главе 5 уточняются и подкрепляются ссылками на литературу основные идеи \textit{Феноменологии внимания} (\textit{Phänomenologie der Aufmerksamkeit}, 2004). Внимание полностью подчинено двойному ритму становления очевидным и становления притягательным, двойное движению пафоса и ответа потому преобретает форму, пронизывающую весь наш опыт и поведение. Что становится очевидным или мыслимым, что доходит до наших глаз или ушей, до наших органов чувств, что до ходит до нас --- не наше. Аллюзия названия \textit{Пороги смысла} (Sinnesschwellen, 1999) указывает, что вездесущее внимание ведет в область эстетики и особого этоса чувств. В последних двух главах говорится также о чуждых аспектах психоанализа и технологии, что занимало значительное место в \textit{Рваных линиях опыта} и \textit{Феноменологии внимания}. Как изначальный опыт пронизан моментами неизведанного и немотивированного, так и феномены в нем с самого начала отмечены подлинной феноменотехнологией и соматотехникой. Технология очень долго --- вопреки всем сетованиям современности об отчуждении в ней --- рассматривалась как расширение сферы собственности, но теперь все больше становится источником особой отчужденности.

Главу 6 можно читать как краткий обзор, собрание различных аспектов чуждого в кратком, ориентированном на отношения между культурами, изложении. Межкультурность не ограничена философией.\footnote{Интересный взгляд на мотив чуждого в современной литературе и теории литературы см. в \textit{Mexiko jenseits der Einsamkeit—Versuch einer interkulturellen Analyse} (1994).} Глобальное мышление не следует ни ожидать, ни желать --- оно неизбежно сформировало бы лишь новые гегемонистские отношения, разве что другими средствами. Философия, переставшая вращаться вокруг себя самой и пустившая корни в привычной культуре, дала бы облегчение. Попытка пересечь границы, не уничтожив их --- одно из будущих начинаний в области межкультурного взаимодействия, и оно значит нечто большее, чем просто сопоставление или смешение культур. Разнообразие слишком легко становится накоплением --- спокойным перебором любопытных вариаций. Опасность для культур, достигних зрелости, заключается в обращении к обществу охотников и собирателей.

Настоящая работа основана на серии отдельных исследований, появившихся в связи с разными обстоятельствами. Все эти исследования были переработаны и ориентированы на общую тему.\footnote{Эти идеи рассматриваются в \textit{Границы нормализации} (\textit{Grenzen der Normalisierung}, 1998).} Я постарался избежать ретроактивной систематизации и сохранить эвристическую множественность подходов, в то же время создав рельефы, благодаря которым --- как на объемной открытке --- были бы видны все значимые аспекты. Как писатели, мы иногда сожалеем, что письмо ограничено двухмерной плоскостью, но, к счастью, в письменной речи --- как и в чтении --- есть свои способы расставления акцентов, и эти подъемы и спуски не дают определенному направлению письма свестись к написанному.

\end{document}
